% Condensed from numerical/tables/ising_small_parameters.tex; numerical values copied without refitting.
\begin{table}[htbp]\centering\footnotesize
\setlength{\tabcolsep}{3.5pt}
\begin{tabular}{lrrrrrrr}\toprule
Direction & $\widehat\alpha$ & $C_{\rm free}$ & $c_0^{\rm CFT}$ & $c_0(F_1)$ & $c_0(F_2)$ & $c_1(F_2)$ & $\eta_0(F_2)$\\\midrule
$I\mid \sigma$ & 2.0001 & 0.20585 & 0.20562 & 0.20573 & 0.20571 & 0.083833 & 0.044771\\
$\sigma\mid I$ & 2.0001 & 0.20582 & 0.20562 & 0.20572 & 0.2057 & 0.069099 & 0.042877\\
$I\mid \varepsilon$ & 3.9934 & 12.491 & 12.988 & 12.83 & 12.883 & -170.71 & -0.80492\\
$\varepsilon\mid I$ & 3.9861 & 12.07 & 12.988 & 12.77 & 12.882 & -355.38 & -0.81263\\
$\sigma\mid \varepsilon$ & 2.0238 & 0.23069 & 0.20562 & 0.2089 & 0.20565 & 12.604 & 0.017273\\
$\varepsilon\mid \sigma$ & 2.0176 & 0.22393 & 0.20562 & 0.20806 & 0.20566 & 9.3133 & 0.022414\\
\bottomrule\end{tabular}
\caption{Ising small-interval fits on $0<x\leq0.02$, with 1255 geometries per direction. The free fit is $C_{\rm free}x^{\widehat\alpha}$; $F_1,F_2$ impose the powers in Table~\ref{tab:num-powers}. All amplitudes include powers of $\pi$. The final column is the signed leading-amplitude discrepancy of $F_2$ in percent, Eq.~\eqref{eq:num-residuals}. The $c_1$ column gives a fitted next amplitude, not a theoretical prediction.}
\label{tab:num-ising-small}
\end{table}
